\documentclass[11pt]{article}

\usepackage[a4paper,margin=1in]{geometry}
\usepackage{amsmath,amssymb}
\usepackage{hyperref}
\usepackage{setspace}

\setstretch{1.12}
\hypersetup{colorlinks=true,citecolor=blue,linkcolor=blue,urlcolor=blue}
\urlstyle{same}

\title{\textbf{BABAPPAi: a likelihood-free framework for diagnosing branch--site episodic-selection identifiability}}

\author{
Krishnendu Sinha\\
Assistant Professor, Department of Zoology, Jhargram Raj College, Jhargram 721507, India
}

\date{}

\begin{document}
\maketitle

\begin{abstract}
Branch--site methods are widely used to infer episodic positive selection from coding-sequence alignments, but their conclusions depend on explicit codon-model assumptions and on whether the available data contain enough information to resolve lineage-restricted signal. Recombination, alignment uncertainty and heterogeneous background constraint can therefore reduce recoverability even when episodic perturbation is biologically present. Here I present BABAPPAi, a likelihood-free diagnostic framework that targets branch--site \emph{identifiability} rather than binary detection. The method combines forward simulation under codon-level mutation--selection dynamics with amortized neural inference trained to recognize branch--site deviations from a neutral or weakly structured background, without gene-wise likelihood fitting or direct estimation of $\omega$. From the inferred branch--site scores, BABAPPAi derives aggregate burden summaries and an Episodic Identifiability Index (EII) that measures whether an alignment contains enough structured variation for episodic interpretation to be statistically meaningful. In held-out synthetic benchmarks, predicted burden increased monotonically with simulated episodic load, remained unstructured under neutrality, and was more stable at the gene level than at individual sites. Performance degraded gradually under recombination-aware segmentation, heterogeneous mutation rates and alignment noise, while deterministic execution and explicit input validation supported robust software behaviour. BABAPPAi is intended as a complement to classical branch--site tests, not a replacement: it diagnoses whether episodic inference is well posed, and at what scale, before strong adaptive claims are made.
\end{abstract}

\noindent\textbf{Keywords:} episodic selection; branch--site evolution; likelihood-free inference; simulation-based inference; mutation--selection models; phylogenomics

\section*{Key points}
\begin{itemize}
    \item BABAPPAi targets the \emph{identifiability} of branch--site episodic signal rather than direct hypothesis testing for positive selection.
    \item The framework combines forward simulation under mutation--selection dynamics with likelihood-free neural inference trained on exact latent labels.
    \item Aggregate burden summaries are more recoverable than site-wise signals, indicating that episodic structure may become measurable at gene scale before it is precisely localizable.
    \item The software reports both $\mathrm{EII}_{z}$ and $\mathrm{EII}_{01}=\sigma(\mathrm{EII}_{z})$, together with \texttt{identifiable\_bool} and \texttt{identifiability\_extent}.
\end{itemize}

\section{Introduction}

Detecting adaptive evolution in protein-coding genes is a central objective of comparative genomics. Codon-based likelihood methods, especially implementations in PAML and HyPhy, have dominated this area because they provide explicit statistical models for nonsynonymous and synonymous substitution processes \cite{Yang2007,Murrell2012,KosakovskyPond2020}. Within this framework, branch--site models are particularly attractive because they allow elevated nonsynonymous rates at a subset of codons along prespecified lineages, matching the biological expectation that adaptation is often sparse in both time and sequence space \cite{AnisimovaYang2007,YangDosReis2011,KosakovskyPond2011}.

Yet classical branch--site inference addresses a narrower target than is often assumed. In practice, these methods do not simply ask whether episodic selection exists in some biological sense; they ask whether the data support that conclusion under a specific codon-model parameterization. This distinction matters because limited alignment length, shallow or excessive divergence, recombination, and alignment error can all degrade statistical recoverability. Earlier simulation studies showed that the performance of likelihood-ratio procedures depends strongly on sequence divergence, model misspecification and alignment quality \cite{Anisimova2001,Wong2004,FletcherYang2010,JordanGoldman2012}. Recombination is especially problematic because most codon-model analyses assume a single genealogy across the alignment, whereas real coding sequences may be mosaics of distinct local histories \cite{Anisimova2003}. As a result, the absence of statistical support is not equivalent to the absence of biological signal.

This gap motivates a different inferential question: \emph{is episodic signal recoverable from the alignment at all, and at what scale?} Recent advances in simulation-based and likelihood-free inference make this question tractable. Neural networks have already shown substantial utility in population-genetic inference because they can learn informative mappings from complex sequence patterns without requiring closed-form likelihoods \cite{SheehanSong2016,Flagel2019,Cranmer2020}. In evolutionary genomics, this is especially useful when the simulator can encode richer biological structure than the downstream analytic model.

BABAPPAi was developed from this perspective. Instead of treating episodic selection as a shift in codon-model $\omega$ classes, the framework represents it as branch- and site-specific perturbation of latent codon-fitness landscapes within a mutation--selection process \cite{HalpernBruno1998}. A neural model is then trained on simulator-generated alignments with exact latent labels to infer branch--site deviation scores from observed sequence patterns. These scores are aggregated into burden summaries and an Episodic Identifiability Index (EII), which measures whether the alignment contains enough structured variation for episodic interpretation to be statistically meaningful. BABAPPAi therefore addresses a complementary problem to classical branch--site testing: not whether a codon-model null can be rejected, but whether episodic signal is recoverable from the data, and at what scale. BABAPPAi is the renamed continuation of the BABAPPA$\Omega$ codebase.

\section{Materials and methods}

\subsection{Framework overview}

BABAPPAi consists of three components: (i) a forward simulator of coding-sequence evolution under branch-specific mutation--selection dynamics, (ii) a neural inference model trained on simulator-generated alignments with exact latent labels, and (iii) a calibration layer that summarizes branch--site scores into aggregate burden measures and the Episodic Identifiability Index. The framework is likelihood-free at inference time: it does not perform gene-wise likelihood optimization, branch-wise $\omega$ estimation, or likelihood-ratio testing.

\subsection{Forward simulation under mutation--selection dynamics}

Evolution is simulated on the 61-sense-codon state space under the standard genetic code, excluding stops. At each site $i \in \{1,\ldots,L\}$, a latent codon-fitness landscape is defined as
\[
F_i:\mathcal{C}\rightarrow \mathbb{R},
\]
where $F_i(c)$ denotes the fitness of codon $c$ at site $i$. Baseline landscapes are sampled as
\[
F_i^{(0)}(c)\sim \mathcal{N}(0,\sigma_F^2),
\]
thereby allowing a continuum of weakly deleterious, nearly neutral and mildly favourable states without imposing discrete codon-model classes.

To introduce gene-wide codon-usage heterogeneity, a codon-specific offset is added:
\[
F_i(c)=F_i^{(0)}(c)+\gamma \log\!\left(\frac{p_c}{\bar p}\right),
\]
where $p_c$ is sampled from a Dirichlet distribution over the codon state space, $\bar p$ is the mean codon frequency, and $\gamma$ modulates the contribution of codon-usage preference.

Episodic selection is generated as a sparse, branch- and site-specific perturbation of the latent fitness landscape. For selected branch--site pairs $(i,b)$,
\[
F_{i,b}(c)=F_i(c)+\Delta F_{i,b}(c),
\qquad
\Delta F_{i,b}(c)\sim \mathcal{N}(0,\sigma_{\mathrm{epi}}^2).
\]
Only a sparse subset of branch--site pairs receives perturbation, preserving the expectation that episodic adaptation is localized in both lineage and position. To approximate local genealogical heterogeneity, the alignment can be partitioned into contiguous blocks that act as coarse recombination-aware segments.

Single-nucleotide mutations follow an HKY-type scheme with transition--transversion bias \cite{Hasegawa1985}. If codons $c$ and $c'$ differ by exactly one nucleotide,
\[
\mu_{c\rightarrow c'}=
\begin{cases}
\kappa\mu, & \text{if the change is a transition},\\
\mu, & \text{if the change is a transversion},
\end{cases}
\]
where $\mu$ is the baseline mutation rate and $\kappa$ is the transition bias parameter. For branch $b$ and site $i$, the selection coefficient of a proposed change is
\[
s_{i,b}(c\rightarrow c')=F_{i,b}(c')-F_{i,b}(c).
\]
Substitutions are then realized under a Halpern--Bruno-style mutation--selection process \cite{HalpernBruno1998}, with continuous-time trajectories simulated using a Gillespie algorithm \cite{Gillespie1977}.

Because episodic perturbations are specified explicitly, exact latent labels are available. For each site--branch pair $(i,b)$, the simulator records an episodic indicator $Y_{i,b}\in\{0,1\}$. These indicators are aggregated into site- and branch-level burdens:
\[
B_i^{\mathrm{site}}=\frac{1}{|\mathcal{B}|}\sum_{b\in\mathcal{B}} Y_{i,b},
\qquad
B_b^{\mathrm{branch}}=\frac{1}{L}\sum_{i=1}^{L} Y_{i,b},
\]
where $\mathcal{B}$ is the set of branches in the tree. These burden quantities provide ground truth for training and benchmarking.

\subsection{Neural inference}

Each training example consists of a codon alignment together with branch-conditioned features derived from parent--child codon transitions, branch lengths and positional context. The model is trained in two stages. First, it is pretrained on a large corpus of simulated alignments spanning neutral, weakly structured and strongly episodic regimes. Second, it is fine-tuned against aggregate burden targets to align network outputs with site- and branch-level latent burden.

Let $\ell_{i,b}$ denote the site--branch logit predicted by the model, and let $\sigma(\cdot)$ be the logistic function. Site- and branch-level predicted burdens are
\[
\hat B_i^{\mathrm{site}}=\frac{1}{|\mathcal{B}|}\sum_{b\in\mathcal{B}}\sigma(\ell_{i,b}),
\qquad
\hat B_b^{\mathrm{branch}}=\frac{1}{L}\sum_{i=1}^{L}\sigma(\ell_{i,b}).
\]
Training minimizes a composite burden-aligned loss,
\[
\mathcal{L}_{\mathrm{total}}=\lambda_s\mathcal{L}_{\mathrm{site}}+\lambda_b\mathcal{L}_{\mathrm{branch}},
\]
where $\lambda_s$ and $\lambda_b$ are fixed loss weights. After training, model parameters are frozen and used as an immutable inference artifact.

At inference time, BABAPPAi returns branch--site deviation scores, site-level burdens, branch-level burdens, and gene-level summaries obtained by averaging site scores within each coding region. These outputs are diagnostic and relative; they are not posterior probabilities of positive selection and they do not replace formal codon-model likelihood statistics.

\subsection{Episodic Identifiability Index}

The framework aims to distinguish the existence of latent episodic perturbation in the generative process from its recoverability in observed sequence data. To summarize this recoverability, BABAPPAi computes an Episodic Identifiability Index (EII).

Let $D_{\mathrm{obs}}$ denote a dispersion statistic derived from the model outputs for a gene, such as the variance of site-level branch-aware logits or an equivalent summary of branch--site heterogeneity. Let $\mu_0$ and $\sigma_0$ be the mean and standard deviation of the same statistic under neutral or near-neutral Monte Carlo replicates matched to the alignment configuration. The standardized effect-size form is
\[
\mathrm{EII}_{z}=\frac{D_{\mathrm{obs}}-\mu_0}{\sigma_0}.
\]

For bounded user-facing reporting, BABAPPAi uses a deterministic monotone transform:
\[
\mathrm{EII}_{01}=\sigma\!\left(\mathrm{EII}_{z}\right)=\frac{1}{1+\exp\!\left(-\mathrm{EII}_{z}\right)}.
\]
By construction, $\mathrm{EII}_{01}\in[0,1]$ and ranking by $\mathrm{EII}_{z}$ and $\mathrm{EII}_{01}$ is identical.

Operational runtime regimes are defined as:
\[
\begin{array}{ll}
0.00 \le \mathrm{EII}_{01}<0.30 & \texttt{not\_identifiable},\\
0.30 \le \mathrm{EII}_{01}<0.70 & \texttt{weak\_or\_ambiguous},\\
0.70 \le \mathrm{EII}_{01}<0.90 & \texttt{identifiable},\\
0.90 \le \mathrm{EII}_{01}\le 1.00 & \texttt{strongly\_identifiable}.
\end{array}
\]
The binary decision field is
\[
\texttt{identifiable\_bool}=\left(\mathrm{EII}_{01}\ge 0.70\right),
\]
and the categorical field \texttt{identifiability\_extent} stores one of the four labels above. These quantities are diagnostic and should not be interpreted as dN/dS estimates, likelihood-ratio statistics, or proof of adaptive substitution.

\subsection{Validation and benchmarking}

Validation was performed on held-out simulated datasets generated with independent random seeds and tree realizations. No alignment used for validation or benchmarking was seen during training or fine-tuning, and all results were obtained with the frozen model in evaluation mode.

Because episodic burden is fundamentally relative, primary evaluation used rank-based criteria. For each alignment, Spearman's rank correlation was computed between true and predicted site burdens and, separately, between true and predicted branch burdens. When the true burden vector was constant, as under strict neutrality, rank correlation is mathematically undefined; such cases were therefore treated as correctly uninformative rather than as failures.

Additional diagnostics included enrichment of true signal among the highest-scoring sites, bootstrap intervals for rank summaries, and a channel-collapse index assessing whether branch-specific outputs degenerated into nearly identical internal representations. To evaluate robustness, stress tests introduced recombination-aware segmentation, heterogeneous mutation rates, alignment noise and malformed input files.

For multi-gene synthetic benchmarks, independently simulated genes were concatenated into larger datasets with known gene-level burdens:
\[
B_g^{\mathrm{true}}=\frac{1}{L_g}\sum_{i=1}^{L_g}Y_i,
\qquad
B_g^{\mathrm{pred}}=\frac{1}{L_g}\sum_{i=1}^{L_g}\hat s_i,
\]
where $L_g$ is gene length, $Y_i$ records latent episodic membership, and $\hat s_i$ is the inferred site score. Four broad burden regimes were considered---neutral, low, medium and high---to assess whether aggregation improved recoverability.

\section{Results}

\subsection{Gene-level burden is more recoverable than site-level burden}

Across held-out non-neutral simulations, BABAPPAi recovered latent episodic burden monotonically as the fraction of perturbed sites increased. The strongest and most stable agreement with ground truth emerged at the gene level, where predicted burden ranked genes according to their simulated episodic load. Site-level resolution was consistently weaker, especially in low-load regimes. This scale dependence is central to the framework: episodic structure may become statistically measurable after aggregation even when individual codons do not yet carry sufficiently stable signal for precise localization.

\subsection{Neutral alignments are correctly uninformative}

Under strictly neutral simulations, true burden is constant by construction and rank-based association is undefined because no ground-truth ordering exists. In this setting, BABAPPAi produced low-dispersion, unstructured outputs rather than spurious burden rankings. This behavior is important because a useful diagnostic should remain uninformative when the data contain no episodic structure to recover.

\subsection{Robustness under biologically plausible model violations}

When recombination-aware block structure, heterogeneous mutation rates and alignment noise were introduced, performance declined relative to the cleanest benchmarks but deteriorated gradually rather than catastrophically. Gene-level ranking remained informative over a broad range of perturbed conditions, whereas fine-scale site localization was more sensitive. This pattern is consistent with previous observations that recombination and alignment error can distort site-wise selection inference in codon-model frameworks \cite{Anisimova2003,FletcherYang2010,JordanGoldman2012}. The results therefore support the use of aggregate diagnostics before strong site-specific claims are made.

\subsection{EII tracks episodic recoverability}

The EII summary increased with simulated episodic burden and remained near the neutral baseline in negative-control datasets. High-EII genes were those in which inferred branch--site dispersion exceeded the calibrated background strongly enough to support meaningful episodic interpretation, whereas low or negative EII values flagged cases in which the data did not justify strong branch--site conclusions. EII therefore functions as a diagnostic gatekeeper: it does not establish adaptation, but it indicates whether the underlying alignment supports statistically interpretable episodic inference.

\subsection{Software behavior is deterministic and robust}

Repeated runs on identical inputs produced identical outputs, indicating deterministic inference after model freezing. Invalid inputs, including empty alignments, non-codon-length sequences and single-taxon files, were rejected with explicit error messages rather than silent corruption. Conserved negative controls generated nearly flat score profiles, whereas synthetic alignments containing localized branch-specific perturbation produced corresponding peaks in inferred burden. These checks support the practical robustness of the implementation.

\section{Discussion}

BABAPPAi is built around a simple but important shift in target: from \emph{detection} of positive selection under a codon-model null to \emph{identifiability} of episodic branch--site structure from the available alignment. These goals are related but not identical. In realistic settings, biological episodic perturbation may be present while the statistical information needed to localize or test it remains weak. A method that explicitly diagnoses this limitation can therefore add value even when classical branch--site tools remain the primary inferential engines.

A second central result is that aggregation stabilizes signal. In the present benchmarks, episodic structure was more recoverable at the gene level than at individual sites. This is both biologically and statistically plausible: episodic adaptation is often sparse, lineage-restricted and transient, so local evidence may be weak even when the gene as a whole carries a real episodic burden. BABAPPAi does not eliminate this difficulty; rather, it quantifies where along the scale hierarchy recoverability begins to emerge.

The EII statistic is useful precisely because it formalizes this question. A high EII indicates that branch--site dispersion exceeds matched neutral expectation strongly enough for episodic interpretation to be statistically meaningful. A low EII indicates that the available data do not support confident episodic interpretation under the chosen calibration. Importantly, EII is not a direct measure of adaptation, nor should negative values be overinterpreted as evidence of purifying selection. Its role is diagnostic: it evaluates whether the branch--site question is well posed for the alignment under study.

The method is intended to complement, not replace, established branch--site procedures such as PAML and HyPhy \cite{Yang2007,Murrell2012,KosakovskyPond2011,KosakovskyPond2020}. In practice, classical likelihood-based tests can still be used to estimate selection effects under explicit codon models, but their outputs may be interpreted more cautiously when BABAPPAi reports low identifiability. Conversely, convergence between strong classical evidence and high EII should increase confidence that the data contain recoverable episodic structure rather than unstable model-driven signal.

Several limitations should be noted. First, the realism of inference depends on the realism of the simulator. Although the simulator incorporates mutation bias, codon-usage heterogeneity, branch-specific perturbation and local segmentation, it cannot capture every feature of real coding-sequence evolution. Second, the current framework yields diagnostic scores rather than formal posterior probabilities or hypothesis tests. Third, the numerical behavior of EII depends on the selected dispersion statistic and on the quality of the matched neutral calibration. Finally, the present manuscript is based on synthetic benchmarking; empirical application will require careful mapping of operational branch labels back onto biologically motivated phylogenetic hypotheses.

\section{Conclusion}

BABAPPAi reframes branch--site analysis around the recoverability of episodic signal. By combining forward simulation under mutation--selection dynamics with likelihood-free neural inference, it diagnoses whether coding-sequence alignments contain enough structured variation for episodic interpretation and identifies the scale at which that structure becomes measurable. The framework does not substitute for classical branch--site testing, but it provides a useful diagnostic layer that helps distinguish biologically absent signal from statistically unrecoverable signal.

\section*{Availability and implementation}

BABAPPAi is available as an open-source Python package at \url{https://github.com/krishnendusinha/babappai} and as the installable package \texttt{babappai} with CLI command \texttt{babappai}. BABAPPAi is the renamed continuation of the BABAPPA$\Omega$ codebase. Legacy frozen-model assets currently remain under the BABAPPA$\Omega$ naming for provenance and reproducibility; the frozen reference model used here is archived at \url{https://doi.org/10.5281/zenodo.18195869}. Legacy software-release provenance is archived at \url{https://doi.org/10.5281/zenodo.18520163}. Benchmarking artifacts, execution logs and related reproducibility materials are archived at \url{https://doi.org/10.5281/zenodo.18197957}.

\section*{Data availability}

All benchmark data used in this study were generated by simulation. The source code, frozen inference artifacts and archived benchmark outputs required to reproduce the analyses are available through the repository and Zenodo records listed above.

\section*{Supplementary data}

Supplementary data, including additional benchmark summaries, stress-test logs and implementation details, should accompany the online version of the manuscript.

\section*{Acknowledgements}

I acknowledge the open-source scientific software ecosystem on which this work depends. I also acknowledge the personal inspiration behind the name ``BABAPPA,'' which originated from my son's imaginative word for a butterfly.

\section*{Funding}

No specific external funding was received for this work.

\section*{Conflict of interest}

The author declares no conflict of interest.

\begin{thebibliography}{99}

\bibitem{Anisimova2001}
Anisimova M, Bielawski JP, Yang Z.
Accuracy and power of the likelihood ratio test in detecting adaptive molecular evolution.
\emph{Mol Biol Evol} 2001;18:1585--1592.

\bibitem{Anisimova2003}
Anisimova M, Nielsen R, Yang Z.
Effect of recombination on the accuracy of the likelihood method for detecting positive selection at amino acid sites.
\emph{Genetics} 2003;164:1229--1236.

\bibitem{AnisimovaYang2007}
Anisimova M, Yang Z.
Multiple hypothesis testing to detect lineages under positive selection that affects only a few sites.
\emph{Mol Biol Evol} 2007;24:1219--1228.

\bibitem{Cranmer2020}
Cranmer K, Brehmer J, Louppe G.
The frontier of simulation-based inference.
\emph{Proc Natl Acad Sci USA} 2020;117:30055--30062.

\bibitem{Flagel2019}
Flagel L, Brandvain Y, Schrider DR.
The unreasonable effectiveness of convolutional neural networks in population genetic inference.
\emph{Mol Biol Evol} 2019;36:220--238.

\bibitem{FletcherYang2010}
Fletcher W, Yang Z.
The effect of insertions, deletions, and alignment errors on the branch-site test of positive selection.
\emph{Mol Biol Evol} 2010;27:2257--2267.

\bibitem{Gillespie1977}
Gillespie DT.
Exact stochastic simulation of coupled chemical reactions.
\emph{J Phys Chem} 1977;81:2340--2361.

\bibitem{HalpernBruno1998}
Halpern AL, Bruno WJ.
Evolutionary distances for protein-coding sequences: modeling site-specific residue frequencies.
\emph{Mol Biol Evol} 1998;15:910--917.

\bibitem{Hasegawa1985}
Hasegawa M, Kishino H, Yano T.
Dating the human-ape splitting by a molecular clock of mitochondrial DNA.
\emph{J Mol Evol} 1985;22:160--174.

\bibitem{JordanGoldman2012}
Jordan G, Goldman N.
The effects of alignment error and alignment filtering on the sitewise detection of positive selection.
\emph{Mol Biol Evol} 2012;29:1125--1139.

\bibitem{KosakovskyPond2011}
Kosakovsky Pond SL, Frost SDW, Muse SV.
A random effects branch-site model for detecting episodic diversifying selection.
\emph{Mol Biol Evol} 2011;28:3033--3043.

\bibitem{KosakovskyPond2020}
Kosakovsky Pond SL, Poon AFY, Velazquez R, Weaver S, Hepler NL, Murrell B, Shank SD, Magalis BR, Bouvier D, Nekrutenko A \emph{et al.}
HyPhy 2.5---a customizable platform for evolutionary hypothesis testing using phylogenies.
\emph{Mol Biol Evol} 2020;37:295--299.

\bibitem{Murrell2012}
Murrell B, Wertheim JO, Moola S, Weighill T, Scheffler K, Kosakovsky Pond SL.
Detecting individual sites subject to episodic diversifying selection.
\emph{PLoS Genet} 2012;8:e1002764.

\bibitem{SheehanSong2016}
Sheehan S, Song YS.
Deep learning for population genetic inference.
\emph{PLoS Comput Biol} 2016;12:e1004845.

\bibitem{Wong2004}
Wong WSW, Yang Z, Goldman N, Nielsen R.
Accuracy and power of statistical methods for detecting adaptive evolution in protein coding sequences and for identifying positively selected sites.
\emph{Genetics} 2004;168:1041--1051.

\bibitem{Yang2007}
Yang Z.
PAML 4: phylogenetic analysis by maximum likelihood.
\emph{Mol Biol Evol} 2007;24:1586--1591.

\bibitem{YangDosReis2011}
Yang Z, dos Reis M.
Statistical properties of the branch-site test of positive selection.
\emph{Mol Biol Evol} 2011;28:1217--1228.

\end{thebibliography}

\end{document}
